\section{Discussion}\label{sec:discussion}

\paragraph{What transfers.}
The portable output of this study is the evidence standard, not the effect
sizes. The unit of claim -- a paired across-seed bits-per-byte delta with a
95\% confidence interval on two corpora, under a single-variable diff at
matched train FLOPs, scored by a versioned suite (Section~\ref{sec:method})
-- is independent of scale and cheap relative to the training runs it
disciplines. Two empirical regularities also seem worth carrying forward.
First, bundle claims should be distrusted by default. The modernized bundle
\citep{grigorev2026imu} supplied this study's motivating headline, a 17.9\%
in-loop validation-perplexity gap at matched 1.19B-token compute (single
runs, in-loop, not suite-stamped; Table~\ref{tab:arms}); under
single-variable iso-FLOP re-tests it decomposed into three small,
individually significant architecture effects of +0.0256 to +0.0355 BPB on
wikitext-2 (each with a zero-excluding 95\% CI, 3 seeds, paired, 131M
tokens per cell, text-lm-v2), one large early-training optimizer-speed
effect (+0.4743 BPB, 95\% CI [+0.4435, +0.5052], 3 seeds, paired, 42M
tokens per cell), and two components -- the WSD schedule and z-loss -- that
contributed nothing measurable (Table~\ref{tab:attrib}). None of that
structure was visible in the bundle run. Second, data interventions
dominated at this scale. The largest attributed effect in the study is a
data swap: dclm-edu over FineWeb-Edu at +0.7034 code BPB (95\% CI
[+0.6639, +0.7430], 3 seeds, paired, 131M tokens per cell, text-lm-v2;
Table~\ref{tab:data}), and the mid-training anneal added +0.2716 code BPB
over its iso-token control (95\% CI [+0.2641, +0.2792], 3 seeds, paired,
150.7M tokens per cell, text-lm-v2; Table~\ref{tab:midtrain}). The
attributed architecture effects are an order of magnitude smaller
(Table~\ref{tab:attrib}), the optimizer effect is scoped to optimization
speed rather than converged quality, and the RL stage returned a
pre-registered null: GRPO beat neither the SFT floor nor the random-reward
gate, with training reward flat at $\approx$0.9\% (math-acc-v1, $n{=}1$
seed, directional by design; Table~\ref{tab:rlvr}, Figure~\ref{fig:rlvr}).
That null is consistent with the sharpen-not-expand reading of RLVR and
with documented spurious-reward effects on Qwen-lineage models
\citep{yue2025does, shao2025spurious} -- consistency, not confirmation,
given the single seed.

\paragraph{What does not transfer.}
The effect sizes themselves. Every comparison in this paper lives between
42M and 1.19B training tokens on one 596M-parameter architecture at
$\approx$2 tokens per parameter, far below compute-optimal
\citep{hoffmann2022training}. These are fixed-budget attributions; their
sign and size may change with budget, architecture, or data mixture, and
the paper makes no scale claim (Section~\ref{sec:limitations}).

\paragraph{The cost of rigor.}
Controls and seeds multiply compute. Computed from the recorded per-cell
budgets (Table~\ref{tab:attrib}), the de-confound ran 12 Phase-1 cells at
131,072,000 tokens each (1,572,864,000 tokens) plus nine Phase-2 cells at
the same budget (1,179,648,000 tokens): $\approx$2.75B training tokens,
roughly 2.3$\times$ the single 1.19B-token bundle run whose result they
re-attributed. That multiple is the price of replacing ``the bundle won''
with ``these components won, and these contributed nothing.'' Nulls cost
the same as wins: the SFT and RLVR stages spent full multi-arm budgets to
establish that nothing separable happened (Table~\ref{tab:sft},
Table~\ref{tab:rlvr}).

\paragraph{What an evidence-gated loop buys.}
The gates run even when no one is looking for a specific error. None of
the seven failures catalogued in Section~\ref{sec:failures} was caught by
re-reading prose; all were caught mechanically -- an implausibility check
on a blend scoring roughly twice the bits-per-byte of both its parents
(Table~\ref{tab:data}), a fixed held-out set that collapsed an in-loop
0.68-PPL ``win'' to +0.009 (Figure~\ref{fig:sft-confound}), and
claim-to-file audits that corrected two already-recorded claims (a
data-ratio pairing and a smoke-log misread) before writing. The same
machinery that catches the errors validates the positive results; the
discipline and the findings are one artifact.
